\chapter{Конструкторская часть}\label{ch:-3}
В данном разделе будут приведены схемы алгоритмов нахождения элемента в массиве


\section{Описание алгоритмов}\label{sec:-3}
На вход алгоритмов подаются 2 строки - str1 и str2, на выходе получается искомое расстояние.
Далее приведены схемы рекурсивных и матричных алгоритмов нахождения расстояний Левенштейна и Дамерау-Левенштейна.

На вход алгоритмы получают массив $arr$ и значение $elem$, которое необходимо найти.
На выходе алгоритмы возвращают найденный индекс, а также количество сравнений.
В случае отсутствия заданного элемента в массиве алгоритмы, в качестве найденного индекса, возвращают $-1$.

На схемах \ref{scheme:lsearch} --- \ref{scheme:bsearch} приведены описания алгоритмов нахождения значения в массиве: использующий полный перебор, использующий бинарный поиск.

\clearpage

\begin{figure}[h]
    \centering
    \includegraphics[scale=0.65]{images/s1}
    \caption{Схема алгоритма поиска полным перебором}
    \label{scheme:lsearch}
\end{figure}

\clearpage

\begin{figure}[h]
    \centering
    \includegraphics[scale=0.5]{images/s2}
    \caption{Схема алгоритма поиска бинарным поиском}
    \label{scheme:bsearch}
\end{figure}

\clearpage


\section{Анализ трудоемкости}

\subsection{Линейный поиск}

Линейный поиск проверяет каждый элемент массива последовательно, пока не найдёт искомый элемент или не проверит все элементы массива.

\begin{enumerate}[label=\arabic*.]
    \item В \textbf{лучшем случае}, элемент может быть найден на первом месте массива. В этом случае потребуется только одно сравнение, и временная сложность будет равна $O(1)$.
    \item В \textbf{худшем случае}, элемент либо находится на последней позиции массива, либо отсутствует в массиве. Тогда алгоритм должен проверить все $n$ элементов массива, что приводит к временной сложности $O(n)$.
\end{enumerate}

\subsection{Бинарный поиск}

Бинарный поиск работает по принципу деления массива пополам на каждом шаге. В начале поиска массив имеет $n$ элементов, и после каждой итерации количество элементов, которые нужно проверить, сокращается в два раза. Рассмотрим временную сложность в различных случаях:

\begin{enumerate}[label=\arabic*.]
    \item В \textbf{лучшем случае}, искомый элемент находится точно в середине массива после первой итерации. Это требует лишь одного сравнения, и временная сложность равна $O(1)$.
    \item В \textbf{худшем случае}, поиск продолжается до тех пор, пока размер массива не будет уменьшен до 1 элемента. Количество делений массива пополам можно выразить как $\log_2{n}$, где $n$ — размер массива. Таким образом, время работы в худшем случае равно $O(\log n)$.
\end{enumerate}

\section*{Вывод}
В данном разделе были представлены схемы алгоритмов нахождения элемента в словаре, а также оценены трудоемкости для данных алгоритмов в лучшем и худшем случаях.

